\chapter{String Rewriting Systems}\label{chap:SRS}
\chapterauthor{Dieter Hofbauer and Johannes Waldmann}

\section{Introduction}


This chapter describes the syntax and semantics for \emph{String Rewriting Systems} (SRSs).

String rewriting is replacement of sub-strings.
Each string can be seen as a term.
For the purposes of the Termination Competition,
each SRS $R$ is translated to a TRS $T$,
and the semantics of $R$ is the TRS semantics of $T$.

\textbf{Example.}
The SRS $R=\{aabb \to bbbaaa\}$
is translated to the TRS $T=\{a(a(b(b(x_1))))\to b(b(b(a(a(a(x_1))))))\}$
over signature $\{a/1, b/1\}$ with variable $x_1$.

\section{Syntax}

A SRS  is a TRS (see Section: TRS/Syntax) where
% \begin{itemize}
% \item
the signature contains unary symbols only,
% \item
and in  each rule, the variable in the left-hand side is equal to the variable in the right-hand side.
% \end{itemize}

\textbf{Example.}
The SRS $R=\{aabb \to bbbaaa\}$ is represented by
\begin{center}
\begin{verbatim}
(format TRS) (fun a 1) (fun b 1)
(rule (a (a (b (b x1)))) (b (b (b (a (a (a x1))))))) 
\end{verbatim}
\end{center}

\section{Semantics}

see Section: TRS/Semantics. All concepts (termination, nontermination),
modifications (full rewriting, innermost and outermost rewriting, relative rewriting)
and measures (complexity) do apply.
In previous Termination Competitions, categories Full Rewriting and Relative Rewriting have been used.

